\documentclass[a4paper,11pt]{article}
\usepackage[utf8]{inputenc}
\usepackage[spanish]{babel}
\usepackage{amsmath,amssymb}
\usepackage{geometry}
\usepackage{hyperref}
\usepackage{listings}
\usepackage{xcolor}
\usepackage{tcolorbox}
% Generador de apuntes: integración con la biblioteca oficial WEC.
% La paleta y las macros visuales (titulo, secciongrande, seccion,
% subtituloA/B, resumen, lista, pasos, definicion, nota, tabla...) proceden
% de bibliotecas_latex/cabecera.tex. Aquí se sincronizan los colores
% del resaltado de código y el índice automático clicable con hyperref.
\providecommand{\cmd}[1]{\texttt{\detokenize{#1}}}
\providecommand{\subnota}[2]{\nota{#1}{#2}}
\providecommand{\subnotasin}[2]{\nota{#1}{#2}}
\AtBeginDocument{\spanishdeactivate{<>}}
\providecommand{\subt}[1]{\par\Needspace{4\baselineskip}\vspace{2mm}\noindent{\bfseries\color{primary}#1}\par\vspace{1mm}}
\providecolor{codeForeground}{HTML}{243247}
\providecolor{codeString}{HTML}{8A4B08}
\providecolor{codeComment}{HTML}{58745C}
\providecolor{codeBackground}{HTML}{F1F7FF}
\providecolor{codeKeyword}{HTML}{2457A6}
\providecolor{consoleBackground}{HTML}{F1F4F8}
\newcommand{\actualizarColoresCodigo}{%
  \colorlet{codeBackground}{\coloractual!7!white}%
  \colorlet{codeKeyword}{\coloractual!80!black}%
}
\newcommand{\iniciarPaletaSecciones}{\pickcolor{\coloractual}}
\AddToHook{cmd/titulo/after}{\iniciarPaletaSecciones\actualizarColoresCodigo}
\AddToHook{cmd/secciongrande/after}{\actualizarColoresCodigo}

% --- ÍNDICE AUTOMÁTICO CLICABLE CON HYPERREF Y CONTROL ANTI-HUÉRFANAS ---
% Se envuelven \secciongrande, \seccion, \subtituloA y \subtituloB usando
% \RenewDocumentCommand para mantener compatibilidad total con hyperref,
% evitar que queden solos al final de página (\Needspace) y garantizar
% que el contenido empiece en la misma página que el título (\nopagebreak).
\makeatletter
\let\orig@secciongrande\secciongrande
\RenewDocumentCommand{\secciongrande}{o m}{%
  \par\Needspace{8\baselineskip}%
  \phantomsection
  \IfValueTF{#1}{\orig@secciongrande[#1]{#2}}{\orig@secciongrande[]{#2}}%
  \nopagebreak\vspace{2mm}\nopagebreak
}
\let\orig@seccion\seccion
\RenewDocumentCommand{\seccion}{o m}{%
  \par\Needspace{6\baselineskip}%
  \phantomsection
  \IfValueTF{#1}{\orig@seccion[#1]{#2}}{\orig@seccion[]{#2}}%
  \nopagebreak\vspace{2mm}\nopagebreak
}
\RenewDocumentCommand{\subtituloA}{m}{%
  \par\Needspace{6\baselineskip}%
  \vspace{4mm}%
  \noindent
  \edef\subtituloreal{\subtitulocolor}%
  \raisebox{0.6ex}{\tikz\fill[\subtituloreal] (0,0) circle (2.2pt);}%
  \hspace{3mm}{\headingfontxb\fontsize{13.5}{16.5}\selectfont\color{\subtituloreal}#1}\par
  \nopagebreak\vspace{2mm}\nopagebreak
}
\RenewDocumentCommand{\subtituloB}{m}{%
  \par\Needspace{5\baselineskip}%
  \vspace{4mm}%
  \noindent
  \edef\subtituloreal{\subtitulocolor}%
  {\headingfont\fontsize{13.5}{16.5}\selectfont\color{\subtituloreal}#1}\par
  \nopagebreak\vspace{2mm}\nopagebreak
}
\makeatother

\lstdefinelanguage{JavaScript}{
  sensitive=true,
  morekeywords={const,let,var,of,async,await,undefined,null,true,false,function,return,if,else,for,while,switch,case,break,continue,new,class,extends,import,export,default,try,catch,throw},
  morecomment=[l]{//},
  morecomment=[s]{/*}{*/},
  morestring=[b]",
  morestring=[b]'
}
\ifcsname codigo\endcsname\else
\newtcblisting{codigo}[1][Java]{listing only,breakable,colback=codeBackground,colframe=\coloractual!28!white,boxrule=0pt,arc=3pt,left=6pt,right=6pt,top=4pt,bottom=4pt,listing options={language=#1,basicstyle=\ttfamily\footnotesize\color{codeForeground},keywordstyle=\color{codeKeyword}\bfseries,morekeywords={const,let,var,of,async,await,undefined,null,true,false,console,document,window,localStorage,fetch,Promise,script,html,head,body,noscript,div,span,p,ul,li,input,form,button,type,src,id,class},stringstyle=\color{codeString},commentstyle=\color{codeComment}\itshape,showstringspaces=false,breaklines=true,columns=fullflexible,keepspaces=true,tabsize=4,upquote=true}}
\fi
\ifcsname consola\endcsname\else
\newtcblisting{consola}{listing only,breakable,colback=consoleBackground,colframe=primary!22!white,boxrule=0pt,arc=3pt,left=6pt,right=6pt,top=4pt,bottom=4pt,listing options={basicstyle=\ttfamily\footnotesize\color{codeForeground},breaklines=true,columns=fullflexible,keepspaces=true,upquote=true}}
\fi

% --- Fallbacks para macros oficiales de la biblioteca en entornos mínimos ---
\providecommand{\paletaAzul}{}
\providecommand{\usarfondo}{}
\providecommand{\titulo}[2]{\begin{center}{\LARGE\bfseries #1}\\[3mm]{\large\bfseries #2}\end{center}\vspace{4mm}}
\providecommand{\coloractual}{blue}
\providecommand{\primary}{blue}
\providecommand{\headingfont}{\bfseries}
\providecommand{\headingfontxb}{\bfseries}
\providecommand{\resumen}[1]{\par\vspace{2mm}\noindent\fbox{\parbox{\dimexpr\linewidth-2\fboxsep-2\fboxrule}{\textbf{Resumen:} #1}}\par\vspace{2mm}}
\providecommand{\definicion}[2]{\par\vspace{2mm}\noindent\fbox{\parbox{\dimexpr\linewidth-2\fboxsep-2\fboxrule}{\textbf{#1:}\ #2}}\par\vspace{2mm}}
\providecommand{\definicionpar}[2]{\par\vspace{2mm}\noindent\fbox{\parbox{\dimexpr\linewidth-2\fboxsep-2\fboxrule}{\textbf{#1:}\ #2}}\par\vspace{2mm}}
\providecommand{\nota}[2]{\par\vspace{2mm}\noindent\fbox{\parbox{\dimexpr\linewidth-2\fboxsep-2\fboxrule}{\textbf{#1:}\ #2}}\par\vspace{2mm}}
\providecommand{\alertas}[1]{\par\vspace{2mm}\noindent\fbox{\parbox{\dimexpr\linewidth-2\fboxsep-2\fboxrule}{\textbf{Alerta:}\ #1}}\par\vspace{2mm}}
\providecommand{\rectangulos}[1]{#1}
\providecommand{\lista}[1]{\begin{itemize}#1\end{itemize}}
\providecommand{\pasos}[1]{\begin{enumerate}#1\end{enumerate}}
\providecommand{\tabla}[2][]{#2}
\providecommand{\inputcarpeta}[1]{%
  \immediate\write18{find "#1" -name "*.tex" 2>/dev/null | sort -V | sed 's/.*/\\input{&}/' > "/tmp/in-#1.tex"}%
  \IfFileExists{/tmp/in-#1.tex}{\input{/tmp/in-#1.tex}}{}%
}


\geometry{top=2.5cm, bottom=2.5cm, left=2.5cm, right=2.5cm}

\title{\textbf{JK}\\ \large Tema 1: Tema 1}
\author{Generador Automático de Apuntes}
\date{\today}

\begin{document}
\maketitle
\tableofcontents
\newpage

\secciongrande{Modelos de ejecución de código: Los cuatro campos del diseño web}

\resumen{
  El diseño web se fundamenta en cuatro pilares: visual, tecnológico, de distribución y propósito. La arquitectura moderna se basa en el modelo cliente-servidor. El desarrollo web es un campo multidisciplinar que requiere dominar todas las fases del proyecto, desde la captación de necesidades hasta la presentación final.
}

\subtituloA{Los cuatro campos del diseño web}

\definicion{Visual}{El aspecto visual de una página web, cuyo objetivo primordial es la comunicación y la atracción de clientes, influyendo directamente en la retención del visitante.}
\definicion{Tecnológico}{La elección de tecnologías (ej. estático o dinámico, bases de datos) que define la implementación del sitio web.}
\definicion{Distribución}{Aspecto crítico relacionado con la velocidad y la fiabilidad de la transmisión de datos entre el servidor y el cliente.}
\definicion{Propósito}{El objetivo fundamental del sitio, que puede ser económico (difusión de productos) o informativo (difusión de eventos).}

\subtituloA{Modelo cliente-servidor}

\definicion{Servidor}{Elemento que aloja hardware, software, técnicas incrustadas y lenguajes de programación, aconsejándose soporte multihilos para la implementación web.}
\definicion{Máquina cliente}{Dispositivo que ejecuta programas (navegadores) para solicitar peticiones al servidor, debiendo soportar HTML, CSS, JavaScript y controles ActiveX.}
\definicion{Conectividad de la red}{El medio y las características utilizadas para la transmisión de datos, cuyo objetivo es que el navegador presente la información alojada en el servidor.}

\nota{Para el examen}{
¿Dónde se implementan los programas para realizar peticiones de solicitud en el servidor?

\textbf{Respuesta:} En la máquina servidor. El navegador se encarga de presentar la información, no de alojarla para responder solicitudes.
}

\secciongrande{Lenguajes de Programación en Entorno Cliente: Dos Entornos de Ejecución}

\resumen{
  La interactividad en portales web requiere lenguajes específicos. Se distinguen entornos cliente (ejecutados en el navegador) y entornos servidor (ejecutados en el servidor). La elección del lenguaje depende de las necesidades del proyecto y de la habilidad del desarrollador.
}

\definicion{Entorno Cliente}{Conjunto de lenguajes ejecutados en el dispositivo del usuario, típicamente un navegador web, para interactuar con la interfaz.}
\definicion{Entorno Servidor}{Conjunto de lenguajes ejecutados en un servidor remoto para gestionar lógica de negocio y bases de datos.}

\subtituloA{Diferenciación y elección de lenguajes}

\lista{
  **Entorno Cliente:** Responsable de la interacción directa con el usuario, mejorando la atractividad de la web. |
  **Entorno Servidor:** Responsable de la lógica de negocio, gestión de datos y respuesta a peticiones del cliente.
}

\nota{Para el examen}{
Es crucial evaluar ventajas e inconvenientes de cada lenguaje; la decisión final recae en el desarrollador según su competencia.
}

\secciongrande{Mecanismos de ejecución de código en un navegador web: La Arquitectura del Navegador}

\resumen{
  El proceso de navegación inicia con la solicitud de un recurso (URL) y finaliza con su visualización. La arquitectura de un navegador, aunque varía por implementación, comparte componentes básicos definidos en una arquitectura de referencia. Estos componentes gestionan desde la interfaz de usuario hasta la ejecución de scripts y la persistencia de datos.
}

\definicion{Arquitectura de Referencia}{Conjunto de componentes básicos comunes a la mayoría de los navegadores web, independientemente de su implementación específica.}

\subtituloA{Fases de la Interacción del Usuario}

\lista{
  **Inicio:** El usuario indica la dirección del recurso (URL/URI). |
  **Proceso:** Depende de la configuración (velocidad, fidelidad al contenido, seguridad). |
  **Finalización:** Visualización exitosa del recurso en la interfaz de usuario.
}

\subtituloA{Componentes Clave de la Arquitectura de Referencia}

\lista{
  \textbf{Interfaz de Usuario:} Gestiona barras de herramientas, progreso y preferencias. |
  \textbf{Motor del Buscador:} Carga la URL y soporta navegación básica (anterior/siguiente). |
  \textbf{Motor de Renderizado:} Interpreta HTML, CSS, XML y presenta visualmente el recurso. |
  \textbf{Subsistema de Comunicaciones:} Implementa protocolos de transferencia (HTTP, etc.). |
  \textbf{Intérprete de JavaScript:} Analiza y ejecuta código JS, manejando eventos. |
  *Parser XML/DOM:** Carga la estructura del documento en memoria como árbol (DOM) para acceso ágil. |
  \textbf{Subsistema de Persistencia:} Almacena historial, sesiones, cookies y certificados de seguridad.
}
\subtituloA{Motores de JavaScript}

\lista{
  **V8:** Utilizado por Chrome, Edge, Opera y Brave. Soporte para Node.js. |
  **SpiderMonkey:** Motor de Firefox. |
  **JavaScriptCore:** Utilizado por Safari (iOS/iPadOS).
}

\nota{Para el examen}{
Los motores modernos (V8, SpiderMonkey, JSC) utilizan compilación Just-In-Time (JIT) a código máquina, lo que garantiza una ejecución muy rápida, superando a los intérpretes lentos antiguos. ¡Recordar que Microsoft abandonó Chakra en 2020!
}

\begin{codigo}
// Ejemplo conceptual de interacción
function cargarRecurso(url) {
  // 1. Motor del Buscador solicita al Subsistema de Comunicaciones
  // 2. Subsistema recibe datos
  // 3. Motor de Renderizado parsea HTML y Parser XML crea el DOM
  // 4. Intérprete de JavaScript ejecuta scripts
  console.log(
\end{codigo}



\secciongrande{CAPACIDADES Y LIMITACIONES DE EJECUCIÓN: ¿Por qué JavaScript en el cliente?}

\resumen{
  JavaScript se utiliza en el lado cliente para mejorar la usabilidad, permitiendo interactividad sin acceder al sistema de archivos del servidor. Es crucial asegurar la navegabilidad sin JS y procesar datos recibidos del servidor para actualizar dinámicamente la interfaz.
}

\subtituloA{Ventajas de JavaScript en el Cliente}

\lista{
  **Mejora de Usabilidad:** Proporciona navegación rápida, moderna y dinámica. |
  **Preparación de Datos:** Prepara datos para su envío al servidor sin modificar archivos en el entorno del servidor. |
  **Resiliencia:** El sitio debe ser navegable incluso si JavaScript está deshabilitado. |
  **Actualización Dinámica:** Examina y actualiza secciones de la página con datos obtenidos del lado del servidor. |
  **Funcionalidades Adicionales:** Ofrece feedback al usuario, interfaz amigable, procesamiento de formularios y modificación dinámica de estilos y contenidos.
}

\subtituloA{Limitaciones y Seguridad (Sandbox)}

\lista{ |
  \textbf{Restricción de Archivos:} No puede escribir ni leer archivos del disco del usuario ni del servidor. | | |
  \textbf{Sandbox:} Toda ejecución de JavaScript está aislada; no puede acceder a la configuración del navegador ni a archivos locales. | | |
  \textbf{Política del Mismo Origen:} Impide que un script de un dominio acceda a recursos de otro dominio (ej. `banco.es` no lee `otrositio.com`). Esto genera el error CORS si el servidor no autoriza. | | |
  \textbf{Vulnerabilidad XSS:} El riesgo principal es inyectar código malicioso (ej. \texttt{<script>}) a través de datos no sanitizados proporcionados por el usuario. }


\nota{Para el examen}{
Recuerde: El navegador protege al usuario de su página, pero no protege su página de los datos maliciosos que usted le introduce (riesgo XSS).
}
\secciongrande{Programación de scripts dónde y cómo colocar los scripts}

\resumen{
  Los scripts de javascript se insertan en un documento html utilizando la etiqueta \texttt{<script>}. Pueden ubicarse tanto en el encabezado (entre \texttt{<head>} y \texttt{</head>}) como en el cuerpo (entre \texttt{<body>} y \texttt{</body>}). La etiqueta de apertura debe incluir el atributo \texttt{type="text/javascript"} para indicar el formato del código.
}

\definicion{Etiqueta script}{Etiqueta contenedora html que designa el inicio y el fin de un bloque de código javascript, comenzando con \texttt{<script>} y terminando con \texttt{</script>}.}
\subtituloA{Ubicación de los scripts en html}

\lista{
  \textbf{Encabezado:} Se coloca entre \texttt{<head>} y \texttt{</head>}. |
  \textbf{Cuerpo:} Se coloca entre \texttt{<body>} y \texttt{</body>}.
}

\subtituloA{Sintaxis y Atributos}

\definicion{Atributo \texttt{type="text/javascript"}}{Indica al navegador que el contenido dentro de \texttt{<script>} está formateado como código JavaScript de texto, no como HTML.}

\begin{codigo}
<script type="text/javascript">
// Código JavaScript aquí
</script>
\end{codigo}


\subtituloA{Ejemplo Básico de Ejecución}

\begin{codigo}
document.write("¡Hola Onliner!");
\end{codigo}



\nota{Para el examen}{
Recordar que el punto y coma (\texttt{;}) al final de una instrucción en JavaScript es crucial para indicar la terminación de la línea al intérprete del navegador.
}
\secciongrande{INTEGRACIÓN DEL CÓDIGO CON LAS ETIQUETAS HTML: Integración online vs. archivo externo}

\resumen{
  Los navegadores permiten integrar JavaScript directamente en HTML (online) o mediante un archivo externo (\texttt{.js}). La integración externa es la práctica más recomendada por su modularidad, facilidad de depuración y rapidez de carga. Es crucial manejar correctamente las rutas de los archivos y considerar el soporte para navegadores sin JavaScript con la etiqueta \texttt{<noscript>}.
}

\definicion{Integración Online}{Inclusión del código JavaScript directamente dentro de las etiquetas \texttt{<script>} del documento HTML, utilizando el atributo \texttt{type="text/javascript"}.}
\definicion{Integración Externa}{Referencia a un archivo JavaScript separado (con extensión \texttt{.js}) mediante el atributo \texttt{src} dentro de la etiqueta \texttt{<script>}.}
\subtituloA{Métodos de Integración}

\lista{
  **Código Online:** %%VERBATIM\_BLOCK\_0%% |
  **Código Externo:** %%VERBATIM\_BLOCK\_1%%
}

\subtituloA{Ventajas de la Integración Externa}

\lista{
  **Independencia de componentes:** Separación clara entre estructura (HTML), formato (CSS) y comportamiento (JS) |
  **Reutilización:** Un mismo archivo JS puede ser compartido por múltiples páginas HTML |
  **Mantenimiento:** Facilita la depuración y mejora la claridad del código |
  **Rendimiento:** Aporta modularidad, seguridad y mayor rapidez en la carga web.
}

\subtituloA{Gestión de Rutas de Archivos}

\lista{
  Ruta Relativa: Se define respecto al directorio actual del archivo HTML (ejemplo: \texttt{../js/ejemplo.js}). Ruta Absoluta: Se define desde la raíz del sistema de archivos o unidad. Ruta Remota: Se utiliza un protocolo (ejemplo: \texttt{http://www.dominio.com/ejemplo.js}).
}
\subtituloA{Soporte y Licenciamiento}

\lista{
  **Sin soporte:** Se utiliza la etiqueta \texttt{<noscript>} para informar al usuario si JavaScript no está habilitado. |
  **Licencias:** Se pueden añadir licencias (ej. Creative Commons) para gestionar el uso del código por otros desarrolladores.
}

\nota{Para el examen}{
¿Por qué es preferible la integración externa? Porque promueve la modularidad, facilita la depuración y mejora la velocidad de carga al separar responsabilidades (estructura, estilo, comportamiento).}

\secciongrande{Herramientas de Programación Sobre Clientes Web}

\resumen{
  La elección del editor de código depende de las necesidades del desarrollador, priorizando herramientas de desarrollo sobre editores de contenido. Es crucial utilizar validadores W3C y entender el modelo cliente-servidor para asegurar la calidad y seguridad del código web.
}

\definicion{Editor de Código}{Herramienta especializada para escribir, editar y gestionar código fuente (HTML, JavaScript, CSS), ofreciendo funcionalidades como resaltado de sintaxis y verificación de errores.}

\subtituloA{Pautas para la selección de editores}

\lista{
  **Buscar facilidad para el desarrollo web:** codificación de colores, verificación de sintaxis, inserción de comentarios y soporte FTP |
  **Ejemplos de editores gratuitos:** Notepad++, Eclipse, NetBeans, KompoZer, Bluefish.
}

\subtituloA{Flujo de trabajo y validación}

\lista{
  Codificar (HTML, JS, CSS) y guardar con extensiones correctas. |
  Ejecutar en 2-3 navegadores distintos (Chrome, Firefox, Safari) para comprobación. |
  Validar el código utilizando el validador del consorcio W3C (http://validator.w3.org/).
}

\nota{Para el examen}{
La estandarización W3C es fundamental para asegurar la compatibilidad entre navegadores. La validación ahorra tiempo y garantiza la validez del código según el estándar web.
}

\end{document}